\section{Introducción}

El análisis de estructuras complejas mediante la teoría de grafos ha cobrado gran relevancia en la última década debido al auge de la ciencia de datos y el análisis de redes. Dentro de este ámbito, el Problema del Clique Máximo (Maximum Clique Problem, MCP) se define formalmente como la búsqueda del subgrafo completo de mayor cardinalidad dentro de un grafo no dirigido $G = (V, E)$, donde todos los vértices del subgrafo están mutuamente conectados \cite{pardalos1994maximum}. Más precisamente, un \emph{clique} en $G$ es un subconjunto $C \subseteq V$ tal que $\{u,v\} \in E$ para todo par $u, v \in C$, y el MCP consiste en encontrar el clique $C^*$ de máxima cardinalidad $\omega(G) = |C^*|$.

La relevancia práctica del problema se sustenta en su amplia aplicabilidad en áreas como la bioinformática, la teoría de códigos, el reconocimiento de patrones, el diagnóstico de fallos, la visión por computadora y la economía \cite{singh2014_simple, pardalos1994maximum}. Sin embargo, el MCP pertenece a la familia de problemas NP-duros \cite{szabo2018_different, akbari_polynomial}, lo que implica que el tiempo computacional crece exponencialmente con el tamaño del grafo y que no se conoce un algoritmo que lo resuelva en tiempo polinomial para instancias arbitrarias. Esta intratabilidad justifica el estudio de estrategias eficientes de poda, ramificación y relajación.

\paragraph{Objetivo general.} Formular el MCP como un Programa Lineal Entero (ILP) y resolverlo de manera exacta sobre una instancia representativa del benchmark DIMACS, evaluando además la calidad de la formulación a través de su relajación lineal continua.

\paragraph{Objetivos específicos.}
\begin{itemize}
    \item Implementar el modelo ILP del MCP en Python utilizando el framework de modelado algebraico Pyomo.
    \item Resolver la instancia \textit{keller4} del benchmark DIMACS mediante el solver de alto rendimiento HiGHS, que aplica internamente Ramificación y Acotamiento (Branch \& Bound) con relajación LP, planos de corte y heurísticas.
    \item Calcular y analizar la brecha de integralidad entre la solución entera y la relajación continua, contrastando los resultados con el óptimo conocido reportado en la literatura.
\end{itemize}

\paragraph{Pregunta de investigación.} ¿Es suficiente la formulación ILP estándar del MCP, resuelta mediante un solver MIP moderno, para alcanzar la solución óptima conocida de una instancia densa del benchmark DIMACS en tiempos razonables, y qué calidad ofrece su relajación lineal continua como cota superior?

El resto del documento se organiza como sigue: la Sección \ref{sec:eda} presenta el estado del arte y dos contextos de aplicación; la Sección \ref{sec:met} describe el método y la formulación matemática; finalmente, se presentan los resultados experimentales, la discusión y las conclusiones.
